\chapter{Introduction}

\section{Motivation}
% Por qué importan los enlaces V2V (avisos de seguridad con poco retardo y alta confiabilidad), cómo reutilizan los canales de la red celular y por qué elegir canal y potencia es un problema difícil que se ha abordado con MARL federado Li et al

\section{Problem statement}
% El esquema de Li et al. decide de forma descentralizada pero aprende a través de la BS (FedAvg y recompensa compartida). Preguntas de investigación: PI1, quitar la BS de la agregación de modelos; PI2, usar una recompensa individual por agente.

\section{Objectives}
% Objetivo general (diseñar, implementar y evaluar la arquitectura descentralizada con recompensa individual) y lista de objetivos específicos.

\section{Contributions}
% Las tres contribuciones: implementación verificada y reproducible del esquema de Li et al., agregación descentralizada por gossip entre vecinos y recompensa individual por agente.

\section{Scope}
% Límites del trabajo: todo es simulación, una sola celda con 15 CUE y 5 a 30 pares V2V, simplificaciones frente al artículo original y posiciones fijas dentro de cada episodio.

\section{How to read this document}
% Para quién está escrito el documento, qué capítulo explica los conceptos desde cero y qué ayudas de lectura tiene (resúmenes de ideas clave, explicación de los símbolos después de cada ecuación).

\section{Document structure}
% Un párrafo que resume qué contiene cada capítulo y cada apéndice.
